% PDF p. 68
\section{实双曲空间}
\subsection{洛伦兹内积}
设 $x=(x_0,x_1,\ldots,x_n)\in\R^{n+1}$，$y=(y_0,y_1,\ldots,y_n)\in\R^{n+1}$。$\R^{n+1}$ 上的洛伦兹内积（Lorentzian inner product）定义为
\begin{equation*}
  \langle x,y\rangle=-x_0y_0+x_1y_1+\cdots+x_ny_n.
  \tag{2.33}\label{eq:2.33}
\end{equation*}
赋予 $\R^4$ 由上述洛伦兹内积诱导的度量张量，所得空间通常也记为 $\R^{1,3}$。四维洛伦兹内积用于狭义相对论，其中 $x_1,x_2,x_3$ 可视为空间变量，$x_0$ 为时间变量。

我们将带有这种洛伦兹内积的 $\R^{n+1}$ 记作 $\R^{1,n}$。

\subsection{双曲面模型（洛伦兹模型）}
\textbf{先介绍双叶双曲面。}考虑 $\R^{n+1}$ 中的方程
\begin{equation*}
  \langle x,x\rangle=-x_0^2+x_1^2+\cdots+x_n^2=-1,
  \tag{2.34}\label{eq:2.34}
\end{equation*}
其中 $x=(x_0,x_1,\ldots,x_n)$，$\langle\cdot,\cdot\rangle$ 是洛伦兹内积。式~\eqref{eq:2.34} 中的 $x_0$ 可满足 $x_0>0$ 或 $x_0<0$，因而此式定义了互不连通的两叶。$x_0>0$ 的部分对应上叶，$x_0<0$ 的部分对应下叶。

\textbf{双曲流形}（hyperbolic manifold）记为 $\mathbb H^n$，定义为双叶双曲面的上叶
\[
  \mathbb H^n=\{x\in\R^{n+1}:\langle x,x\rangle=-1,\quad x_0>0\}.
\]
通常取上叶，因为它是连通的，而且可以避免定义双曲空间时的歧义。取 $x_0>0$ 简化了分析，也保证定向一致。

% PDF p. 69
\begin{statement}{命题 2.1}{}
若 $p\in\R^{n+1}$ 且 $\langle p,p\rangle<0$，则洛伦兹内积 $\langle\cdot,\cdot\rangle$ 在正交补 $p^\perp$ 上的限制，即 $\langle p^\perp,p^\perp\rangle$，是正定的。

\textbf{证明。}先给出正交补 $p^\perp$ 的定义。
\[
  p^\perp=\{v\in\R^{n+1}:\langle p,v\rangle=0\},
\]
其中 $\langle\cdot,\cdot\rangle$ 表示洛伦兹内积。目标是证明 $\langle\cdot,\cdot\rangle$ 在 $p^\perp$ 上的限制是\textbf{正定的}，即任取非零 $v\in p^\perp$，有 $\langle v,v\rangle>0$。

取 $v_\perp\in p^\perp$。任意向量 $v\in\R^{n+1}$ 可写成向量 $p$ 和 $v_\perp$ 的组合
\begin{equation*}
  v=\alpha p+v_\perp,
  \tag{2.35}\label{eq:2.35}
\end{equation*}
其中 $\alpha\in\R$ 为标量。由 $\langle p,v_\perp\rangle=0$ 得
\begin{equation*}
  \begin{aligned}
    \langle v,v\rangle
      &=\langle\alpha p+v_\perp,\alpha p+v_\perp\rangle\\
      &=\alpha^2\langle p,p\rangle+\langle v_\perp,v_\perp\rangle.
  \end{aligned}
  \tag{2.36}\label{eq:2.36}
\end{equation*}
若 $v\in p^\perp$，则由式~\eqref{eq:2.35} 得
\begin{equation*}
  \langle p,v\rangle=\alpha^2\langle p,p\rangle.
  \tag{2.37}\label{eq:2.37}
\end{equation*}
但 $\langle p,v\rangle\geq0$，这意味着 $\alpha=0$。因此
\begin{equation*}
  \langle v,v\rangle=\langle v_\perp,v_\perp\rangle.
  \tag{2.38}\label{eq:2.38}
\end{equation*}
回顾
\begin{equation*}
  \langle p,p\rangle=-p_0^2+p_1^2+\cdots+p_n^2.
  \tag{2.39}\label{eq:2.39}
\end{equation*}
于是
\begin{equation*}
  \lim_{p_0\to-\infty}\langle p,p\rangle\longrightarrow-\infty.
  \tag{2.40}\label{eq:2.40}
\end{equation*}
若 $v\in p^\perp$，则 $\langle v,p\rangle=-v_0p_0+v_1p_1+\cdots+v_np_n=0$，从而
\begin{equation*}
  v_0=\frac{v_1p_1+\cdots+v_np_n}{p_0}.
  \tag{2.41}\label{eq:2.41}
\end{equation*}

% PDF p. 70
由式~\eqref{eq:2.41} 得
\begin{equation*}
  \lim_{p_0\to-\infty}v_0=\frac{v_1p_1+\cdots+v_np_n}{p_0}\longrightarrow0.
  \tag{2.42}\label{eq:2.42}
\end{equation*}
式~\eqref{eq:2.42} 意味着
\begin{equation*}
  v\longrightarrow(v_1,\ldots,v_n).
  \tag{2.43}\label{eq:2.43}
\end{equation*}
因此，由式~\eqref{eq:2.38} 与式~\eqref{eq:2.43} 得
\begin{equation*}
  \langle v_\perp,v_\perp\rangle=\langle v,v\rangle=v_1^2+\cdots+v_n^2>0,
  \tag{2.44}\label{eq:2.44}
\end{equation*}
其中 $v\in p^\perp\setminus\{0\}$。这就证明了命题 2.1。\translatornote{式~\eqref{eq:2.37} 右边的系数按双线性应为 $\alpha$。$v\in p^\perp$ 时左边为 $0$。式~\eqref{eq:2.42} 的极限不能推出对固定 $p$ 的式~\eqref{eq:2.44}。对固定 $p$，式~\eqref{eq:2.41} 与柯西–施瓦茨不等式给出 $v_0^2\leq (\sum_{i=1}^np_i^2/p_0^2)\sum_{i=1}^nv_i^2$，而 $\sum_{i=1}^np_i^2/p_0^2<1$。因此非零 $v\in p^\perp$ 确有 $\langle v,v\rangle>0$。}
\end{statement}

由于双曲面在点 $p$ 处的切空间为 $p^\perp$，黎曼度量 $g$ 可由洛伦兹度量定义。因而有以下推论。

\begin{statement}{推论 2.1}{}
设 $u,v\in p^\perp$。用洛伦兹度量定义黎曼度量
\begin{equation*}
  g=\langle u,v\rangle.
  \tag{2.45}\label{eq:2.45}
\end{equation*}
\end{statement}
